% !TeX root = ../../Book3.tex
% 纲要标题：### B.2 高效 Gröbner 算法

\section{高效 Gröbner 算法}
\label{b3:sec:B02}

逐项求余会重复消去许多相同单项式，而系数增长也可能使精确运算变慢。把约化组织成矩阵、转到有限域计算再重构，都需要保留能够回到原系数环核验的证据。

% 纲要细目（与计划纲要同步）：
% * F4、F5 型算法、签名准则及稀疏线性代数
% * FGLM 的零维换序与适当的 Gröbner walk
% * 模方法、坏素数、重构及原系数域验证
% * 实现细节可以替换，数学正确性证书不得依赖单一软件品牌
%
% #### B.2.1 F4、F5 型算法、签名准则与稀疏线性代数
% #### B.2.2 FGLM 零维换序与 Gröbner walk
% #### B.2.3 模方法、坏素数、重构与原域验证
% #### B.2.4 算法实现与结果验证


\subsection{F4、F5 型算法、签名准则与稀疏线性代数}
\label{b3:subsec:B02:01}

Buchberger 算法逐对形成 \(S\)-多项式并求余，其中许多约化会反复使用同一个单项式倍数。F4 型方法把一批这样的运算组织成线性代数：先选取临界对，收集所需多项式倍数；以出现的单项式为列、这些倍数的系数为行；再作行消元，提取产生新首项的行。每个行变换都是系数域上的可逆线性变换，因此保留这批多项式张成的向量空间，也保留它们属于原理想的表示。

\ProseParagraphBreak
这里还有一个不可省略的步骤：收集用于消去首项的约化行。例如一行出现了可被现有首项整除的单项式，就应加入相应的基元素倍数，再检查新引入的项。这称为符号预处理。若只把若干输入多项式随意排成矩阵，行最简形并不会自动成为整个理想的 Gröbner 基。临界对的选择、预处理以及尚未处理的对仍需统一管理。本节说明的是 F4 型方法的运算原理，不把这几句话作为一个完整的 F4 实现或其全部优化的证明。

\begin{definition}\label{b3:def:B-signature}
设 \(k\) 为域，\(S=k[x_1,\ldots,x_n]\)，取 \(f_1,\ldots,f_s\in S\)，并在 \(S^s\) 上指定全局模项序。对于带有表示
\[
(u,f),\qquad u=(u_1,\ldots,u_s),\quad f=\sum_i u_if_i,\quad u\ne0
\]
的数据，称 \(\operatorname{LM}(u)\) 为这份表示的\term{签名}[signature]。同一个多项式的不同表示可以有不同签名。
\end{definition}

例如 \(f_1=x,f_2=y\) 时，\(y e_1-x e_2\) 表示零多项式，却携带一个非零模首项。它说明某些高签名的表示已有关系，可以用更低签名的表示替代。

\begin{proposition}\label{b3:prop:B-signature-rewrite}
设 \(k\) 为域，\(S=k[x_1,\ldots,x_n]\)，取 \(f_1,\ldots,f_s\in S\)，令 \(F:S^s\to S\) 为 \((u_i)\mapsto\sum_i u_if_i\)，并固定 \(S^s\) 上的全局模项序。若非零 \(z\in\ker F\) 的首单项式整除非零向量 \(u\in S^s\) 的首单项式，则 \(Fu\) 有一个向量表示 \(u'\)，满足 \(u'=0\) 或 \(\operatorname{LM}(u')<\operatorname{LM}(u)\)。
\end{proposition}
\begin{proof}
两个首单项式的整除意味着它们在同一基分量中，故可取一个带系数单项式 \(c x^\alpha\)，使 \(u\) 与 \(c x^\alpha z\) 的首项相同。令 \(u'=u-cx^\alpha z\)，首项消去后得到所需严格不等式，而 \(Fu'=Fu\)，因为 \(Fz=0\)。
\end{proof}

F5 型及其他签名算法把这种关系信息用于排除不必再次处理的候选表示，并限制约化步骤，使签名信息能够继续追踪。上面的命题只证明可以降低一份表示的签名，并没有说它所表示的多项式可以无条件丢弃；完整算法还须保证较低签名的代表会被处理。

\ProseParagraphBreak
签名准则的安全性依赖指定的模序、已知合冲关系、重写规则及处理次序。不能仅因两份多项式的某个签名相同，就在没有这些不变量的情况下任意丢弃一份；也不能由算法名称推断任意输入均不会出现约化为零。完整的 F4、F5 算法及相应准则属于进一步的算法专题。对本书已经构造的候选答案，\subsecref{b3:subsec:B02:04}给出的普通 Gröbner 证书仍足以独立核验，不要求验证者复现生成答案时采用的优化。

\subsection{FGLM 零维换序与 Gröbner walk}
\label{b3:subsec:B02:02}

\subsecref{b3:subsec:0805:03}已把换序解释为商代数中的线性相关判定。现在给出完整的有限算法。它不要求理想为根理想，也不要求 \(k\) 代数闭。

\begin{algorithm}[htbp]
\caption{零维商代数的 FGLM 换序}\label{b3:alg:B-fglm}
\begin{algorithmsteps}
{精确可计算域 \(k\)，真理想 \(I\subset S=k[x_1,\ldots,x_n]\) 的一个全局 Gröbner 基，满足 \(D=\dim_k(S/I)<\infty\)，以及目标全局项序 \(<\)。}
{理想 \(I\) 关于 \(<\) 的 Gröbner 基 \(G\)，以及新标准单项式表 \(B\)。}
\State 在旧标准单项式基中计算乘法矩阵 \(T_i\)，记 \(v(m)\in k^D\) 为单项式 \(m\) 的坐标。
\State 置有序候选集合 \(Q=\{1\}\)，已处理集合 \(P=\varnothing\)，并置 \(B=G=\varnothing\)。
\While{\(Q\ne\varnothing\)}
\State 取出 \(Q\) 中最小的单项式 \(m\)，将 \(m\) 记入 \(P\)。
\If{\(m\) 被某个 \(g\in G\) 的首单项式整除}
\State 丢弃 \(m\)，继续取下一个候选。
\Else
\State 计算 \(v(m)\)，判定它与已保存的 \(v(b)\ (b\in B)\) 是否线性相关。
\If{\(v(m)=\sum_{b\in B}c_bv(b)\)}
\State 把 \(m-\sum_{b\in B}c_bb\) 加入 \(G\)。
\Else
\State 把 \(m\) 加入 \(B\)，并把尚未属于 \(Q\cup P\) 的 \(x_1m,\ldots,x_nm\) 加入 \(Q\)。
\EndIf
\EndIf
\EndWhile
\State 返回 \(G,B\)。若需要约化 Gröbner 基，最后作首一化与交互约化。
\end{algorithmsteps}
\end{algorithm}

队列始终按目标序取最小元，但不要求枚举这个序中位于某个单项式之前的全部单项式。每个候选只从已接受的基单项式的相邻倍数产生。坐标可由 \(v(x_i b)=T_i v(b)\) 计算，也可用旧基求余；保存坐标可以避免重复运算。若 \(I=S\)，旧基已检出 \(1\)，直接返回 \(\{1\}\) 和空基即可，不进入上述真理想版本。

\begin{theorem}\label{b3:thm:B-fglm}
设 \(k\) 是支持精确四则运算和零判定的域，\(I\subset S=k[x_1,\ldots,x_n]\) 为真理想，且 \(D=\dim_k(S/I)<\infty\)。给定 \(I\) 关于一个可有效比较的全局项序的 Gröbner 基，并指定另一可有效比较的全局项序 \(<\)。则\cref{b3:alg:B-fglm}终止，返回的 \(B\) 是 \(S/I\) 的目标标准单项式基，\(G\) 是 \(I\) 关于 \(<\) 的 Gröbner 基。
\end{theorem}
\begin{proof}
每次加入 \(B\) 时，坐标向量与原有向量线性无关，故这种加入至多发生 \(D\) 次。除初始候选 \(1\) 外，每个候选由某次加入产生，因而总共至多有 \(1+nD\) 个候选。每个候选只处理一次，所以算法终止。

全局项序满足 \(m<x_i m\)。由队列取最小元和只从已处理单项式产生更大单项式可知，处理次序严格递增。因此写入 \(G\) 的关系 \(g=m-\sum c_bb\) 确以 \(m\) 为首单项式；它的旧坐标为零，故 \(g\in I\)。最终 \(B\) 的元素都不被 \(G\) 的首单项式整除：在它加入时不被已有首项整除，以后产生的首项又比它大，不可能整除它。

反过来，凡不被这些首单项式整除的单项式都在 \(B\) 中。否则取这样的最小单项式 \(m\)。\(m\ne1\)，因为真理想情形下 \(v(1)\ne0\)，于是 \(1\) 被接受。写 \(m=x_i b\)，其中 \(b\) 是一个真因子。\(b\) 也不被任何首单项式整除，且 \(b<m\)，故由最小性 \(b\in B\)。于是 \(m\) 已进入队列。处理它时，既不能按首项整除丢弃，也不能因相关而产生以 \(m\) 为首项的关系，只能加入 \(B\)，矛盾。

令 \(J=(G)\subset I\)。全局多项式除法把每个 \(f\in S\) 写成 \(G\) 的组合加上由 \(B\) 张成的余项。若 \(f\in I\)，这个余项在 \(S/I\) 中为零；而 \(B\) 的坐标线性无关，故余项为零，得到 \(I=J\)。若非零 \(f\in I\) 的首单项式不被任何 \(\operatorname{LM}(g)\) 整除，除法会把它保留为余项的最高项，后续较小项无法消去它，与余项为零矛盾。因此 \(G\) 是 Gröbner 基。由余项的张成性及坐标独立性，\(B\) 是商空间基，也是目标标准单项式集合。
\end{proof}

\begin{example}\label{b3:exm:B-fglm-complete}
取 \(I=(x-y^2,y^3-2)\subset\mathbb Q[x,y]\)，旧序为 \(x>y\) 的字典序，旧基为 \(1,y,y^2\)。在这些坐标中
\[
T_y=\begin{pmatrix}0&0&2\\1&0&0\\0&1&0\end{pmatrix},
\qquad T_x=T_y^2.
\]
改用 \(y>x\) 的字典序。接受 \(1\) 后，候选为 \(x,y\)；先取 \(x\)，其坐标为 \((0,0,1)^{\mathsf T}\)，接受。接着取 \(x^2\)，其坐标为 \((0,2,0)^{\mathsf T}\)，仍接受。再取 \(x^3\)，其坐标为 \((4,0,0)^{\mathsf T}\)，所以产生 \(x^3-4\)，不再从 \(x^3\) 产生后继。

\ProseParagraphBreak
随后取 \(y\)，因 \(v(y)=\tfrac12v(x^2)\)，产生 \(y-\tfrac12x^2\)。已在队列中的 \(xy,x^2y\) 被首项 \(y\) 整除，全部丢弃。最后
\[
B=(1,x,x^2),\qquad
G=\{x^3-4,\ y-\tfrac12x^2\}.
\]
两个关系都可在旧商环直接核验。它们的首项 \(x^3,y\) 给出恰好三个标准单项式，与算法保存的独立坐标一致。
\end{example}

有限维是上述队列计数的根据。对正维理想，标准单项式无限，不能照搬终止证明。Gröbner walk 则采取另一种思路：用权向量及必要的细化次序描述项序，在一段权向量变化中首项不变；跨过变化的位置时，先计算相应初始理想，再把关系提升回原理想。它涉及 Gröbner 扇及权初始理想的完整理论，本附录只说明与 FGLM 的区别，不给出一般 walk 算法，也不把它作为后文证明的依据。

\subsection{模方法、坏素数、重构与原域验证}
\label{b3:subsec:B02:03}

在 \(\mathbb Q[x]\) 中计算时，可以先清除输入分母，在有限域上试算，再尝试恢复有理系数。这样做减少了中间系数膨胀，但必须区分三件事：某个素数下的计算正确，若干素数下的答案一致，以及重构答案在 \(\mathbb Q\) 上正确。

\ProseParagraphBreak
坏素数可能使首系数变零、分母不可逆，或使某个秩下降。例如输入 \(2x+y,y^2-1\) 在特征不为 \(2\) 时给出 \(x=-y/2\)，而模 \(2\) 后成为 \(y,y^2-1\)，生成整个环。这个素数下的正确答案不能描述原有的零维商代数。

\begin{proposition}\label{b3:prop:B-good-primes-certificate}
设 \(F\) 为 \(\mathbb Q[x_1,\ldots,x_n]\) 的一个有限多项式列，固定全局项序，\(G\) 是同一理想的首一 Gröbner 基。固定 \(G\) 由 \(F\) 的表示、\(F\) 由 \(G\) 的表示，以及所有 \(S\)-多项式约化为零的有限除法记录。存在非零整数 \(N\)，使这些记录可在 \(\mathbb Z[1/N]\) 中书写，且对每个 \(p\nmid N\)，模 \(p\) 后的 \(G\) 仍为模 \(p\) 后 \(F\) 的 Gröbner 基，并保留首单项式。
\end{proposition}
\begin{proof}
取 \(N\) 为全部输入、基元素及有限记录所出现系数的分母的公倍数，并包含使各首项标准化所用的非零整数因子。模 \(p\nmid N\) 时，所有等式有意义。首一性使基的首项不消失；两向表示证明两列生成同一理想；约化记录中的单项式上界仍成立，故 Buchberger 判据给出结论。首单项式既未消失也未改变，标准单项式集合随之保持。
\end{proof}

这个命题解释了好素数为何存在，却不是计算之前就已知 \(N\) 的办法：其证明使用了最终答案的证书。实际计算可先按首项形状分组，利用中国剩余定理合并同形答案，再作有理重构。若模 \(M\) 的剩余类 \(c\) 要恢复为 \(a/b\)，可要求
\[
a\equiv cb\pmod M,\quad |a|\le A,\quad 0<b\le B,\quad \gcd(a,b)=1.
\]
当 \(2AB<M\) 且分母可逆时，范围内答案至多一个，因为两个答案的交叉差既被 \(M\) 整除，又绝对值小于 \(M\)。但合理的大小界未必预先可知，重构失败时应继续加入素数。

\ProseParagraphBreak
最终的有理候选仍须回到原系数域核验两向生成关系与 Gröbner 条件。许多素数下结果一致可以帮助选择候选，却不能替代这一步；验证成功后，结论就不依赖猜测哪些素数是好的。

\subsection{算法实现与结果验证}
\label{b3:subsec:B02:04}

设输入列为 \(F=(f_1,\ldots,f_s)\)，候选列为 \(G=(g_1,\ldots,g_t)\)。一份可单独检查的证书应包含工作域、变量和项序，以及下列有限数据：
\begin{enumerate}
\item 一个矩阵 \(U\)，使 \(G=FU\)，证明每个候选仍在输入理想中。
\item 对每个 \(f_i\) 的表示 \(f_i=\sum_jv_{ji}g_j\)，即 \(F=GV\)，证明没有丢掉原理想。
\item 每个所需 \(S\)-多项式的零余项记录，其各乘积首单项式严格小于相应首单项式的最小公倍数。
\end{enumerate}
前两项给出理想相等，第三项由 Buchberger 判据证明 \(G\) 确为该项序的 Gröbner 基。若不采用临界对删减准则，直接检查全部有限多个对即可；若采用删减，也应给出相应判据或补回被省去的验证。

\ProseParagraphBreak
保存表示矩阵并不要求每次重新求解理想成员问题。初始输入附单位向量；每次乘单项式、相减和归一化时，对表示向量作同样运算，即可始终保存 \(g=Fu\)。F4 的行操作可以同样作用于表示矩阵。FGLM 返回的每个关系可先用旧基求余并取得表示，再与旧基到输入的变换矩阵相乘。

\begin{example}
只检查 \(F\) 对 \(G\) 的余项为零是不够的：任意输入对 \(G=\{1\}\) 都有零余项，但输入理想未必为整个环。反过来，只检查 \(G\subset(F)\) 也不够：空集或全零列不会恢复非零输入理想。两个方向的矩阵等式正是排除这两种错误的必要信息。
\end{example}

成员证明则更简单：只要明确写出 \(f=\sum a_if_i\)，即可逐项展开核验，不需要先证明给定生成组为 Gröbner 基。非成员结论需要更多结构，例如对已经验证的 Gröbner 基取得非零正规形。答案的生成方法可以很复杂，而最后检查这些等式的过程可以保持直接。
